\documentclass[10pt]{article}
\usepackage[a4paper,margin=2.5cm]{geometry}
\usepackage[T1]{fontenc}
\usepackage{lmodern,amsmath,amssymb,booktabs}
\begin{document}
\section*{Task 1: T16 reference model B3}
Let $K=\{0,\ldots,40\}$ be the VMs, $L=\{0,\ldots,8\}$ the servers,
and $U=\{0,0.4,0.7,0.9\}$ the tangent points. Data are generated by
\texttt{cp34\_instance(16,'B',3)}. VM $k$ needs CPU $q_k$ and memory $r_k$;
server $l$ has CPU capacity $C_l$ and memory capacity $R_l$.
Let $g_k$ indicate a GPU VM and $G_l$ a GPU server.
The reference model has binary assignments $y_{kl}$ and continuous delays
$t_l\ge0$; the LP replaces binary assignments by $0\le y_{kl}\le1$.
Define $a_u=(1-u)^{-2}$ and $b_u=-u^2(1-u)^{-2}$.
\begin{align*}
\min\quad &\sum_{l\in L}t_l\\
\text{B1: }&\sum_{l\in L}y_{kl}=1 &&(k\in K),\\
\text{B2: }&\sum_{k\in K}q_k y_{kl}\le C_l &&(l\in L),\\
\text{B3: }&\sum_{k\in K}r_k y_{kl}\le R_l &&(l\in L),\\
\text{B4: }&a_u\sum_{k\in K}\frac{q_k}{C_l}y_{kl}-t_l\le-b_u
&&(l\in L,\ u\in U).
\end{align*}
B1 assigns each VM completely in total; B2 and B3 enforce CPU and memory
capacities. B4 bounds latency below by each prescribed tangent to
$g(u)=u/(1-u)$. If $g_k=1$ and $G_l=0$, set the upper bound of $y_{kl}$
to zero; retain the variable. No zone constraints are added to B3.

\paragraph{Matrix form and indexing.}
Use the stacked vector $x=(y_{00},\ldots,y_{0,8},y_{1,0},\ldots,y_{40,8},t_0,\ldots,t_8)$.
Thus $j(y_{kl})=9k+l$, $j(t_l)=369+l$, and
$c=(0_{369},1_9)$. There are 41 equality rows followed by 54
``$\le$'' rows. Row indices are $k$ for B1, $41+l$ for B2,
$50+l$ for B3 and $59+4l+m$ for B4, where $m=0,\ldots,3$
indexes the ordered tangent points.
Hence $A\in\mathbb R^{95\times378}$. Lower bounds are zero;
assignment upper bounds are one, or zero for forbidden pairs;
$t_l$ has no upper bound. Bounds are not counted as rows.
\begin{center}\begin{tabular}{lr}\toprule
Check & Value\\\midrule
Servers / VMs & 9 / 41\\
Total CPU capacity / request & 420 / 201\\
Total memory capacity / request & 1339 / 706\\
Variables / constraint rows & 378 / 95\\
Forbidden assignments retained with upper bound zero & 35\\
B1 / B2 / B3 / B4 rows & 41 / 9 / 9 / 36\\\bottomrule
\end{tabular}\end{center}
The designated sensitivity row is \texttt{assign\_9}, matrix row 9,
with RHS 1; VM 9 requests 8 CPU units. B4 was multiplied by $-1$
relative to its ``$\ge$'' form. Accordingly, the dual price of a B4
matrix row has the opposite sign to that of the original row.

\newpage
\section*{S1 --- The model and what it ignores}
The T16 B3 reference LP has 41 VMs and 9 servers. It contains
$41\cdot9+9=378$ variables and $41+9+9+36=95$ constraint rows,
matching the team check values. All 35 GPU-incompatible assignments
remain in the model with upper bound zero; bounds are not extra rows.
The CPU and memory totals are $201/420$ and $706/1339$ (request/capacity).
These aggregate totals alone do not establish feasibility: individual
capacities and GPU compatibility still restrict the assignments.

The objective minimizes the sum of server latency variables, each bounded
by four prescribed tangents at utilization $0,0.4,0.7,0.9$.
This piecewise-linear lower approximation replaces the exact M/M/1
function $u/(1-u)$ and the squared-utilization objective used in our
Checkpoint 2 CPU block. Its value therefore cannot be directly compared
with the Checkpoint 2 optimum as a measure of improvement.
The LP also permits VM splitting, which is a mathematical relaxation
rather than an executable placement. GPU restrictions and memory capacities
are retained; no zone restrictions are introduced for variant B3.
The formulation is kept unchanged so its LP bound remains comparable
with the integer reference model in Checkpoint 4.
\begin{flushright}\textit{Dilnaz Bekturova}\end{flushright}

\newpage
\section*{H1 --- Model owner hand task}
\textbf{Dilnaz Bekturova, Team T16.}
\paragraph{(a) LP matrix form.}
For $y=(y_1,y_2,y_3)^T$,
\[
\min c^Ty,\qquad Ay\ge b,\qquad 0\le y\le1,
\]
\[
c=\begin{pmatrix}3\\5\\6\end{pmatrix},\quad
A=\begin{pmatrix}1&1&0\\0&1&1\\1&0&1\end{pmatrix},\quad
b=\begin{pmatrix}1\\1\\1\end{pmatrix}.
\]
There are three variables and three constraint rows; variable bounds
are not counted as rows.
\paragraph{(b) Conversion to ``$\le$''.}
Multiplying each row by $-1$ gives
\[
\begin{pmatrix}-1&-1&0\\0&-1&-1\\-1&0&-1\end{pmatrix}y
\le\begin{pmatrix}-1\\-1\\-1\end{pmatrix}.
\]
For minimization, original ``$\ge$'' row prices satisfy $\pi_i\ge0$.
The converted RHS is $\widetilde b=-b$, so
$\widetilde\pi_i=\partial z/\partial\widetilde b_i=-\pi_i\le0$.
For example, $(1,4,2)$ becomes $(-1,-4,-2)$; the reduced-cost
expression is unchanged because $(-A)^T(-\pi)=A^T\pi$.
\paragraph{(c) All binary points.}
Each row below evaluates all three covering sums and the objective.
\begin{center}
\begin{tabular}{ccccc}
\toprule
$y$ & $(y_1+y_2,y_2+y_3,y_1+y_3)$ & Feasible? & $3y_1+5y_2+6y_3$\\
\midrule
$(0,0,0)$ & $(0,0,0)$ & No & 0.000\\
$(0,0,1)$ & $(0,1,1)$ & No & 6.000\\
$(0,1,0)$ & $(1,1,0)$ & No & 5.000\\
$(0,1,1)$ & $(1,2,1)$ & Yes & 11.000\\
$(1,0,0)$ & $(1,0,1)$ & No & 3.000\\
$(1,0,1)$ & $(1,1,2)$ & Yes & 9.000\\
$(1,1,0)$ & $(2,1,1)$ & Yes & 8.000\\
$(1,1,1)$ & $(2,2,2)$ & Yes & 14.000\\
\bottomrule
\end{tabular}
\end{center}
Thus $z^*=\min\{11,9,8,14\}=8.000$, attained at $y=(1,1,0)$.
\begin{flushright}
\textit{Dilnaz Bekturova. Signature: \underline{\hspace{3cm}}.}\par
\textit{Date: \underline{\hspace{3cm}}.}
\end{flushright}
\end{document}
